\documentclass[12pt]{article}
\usepackage[utf8]{inputenc}
\usepackage[T1]{fontenc}
\usepackage{hyperref}
\title{Kui Language Resources for K-12 Curriculum Development}
\author{Compiled by Claude Sonnet 5 (claude-sonnet-5) for Andr\'as Kornai}
\date{}
\begin{document}
\maketitle
\section{Introduction}
Kui (ISO 639-3 uki) is a South-Central Dravidian language of Odisha, India. It is written in the Odia script. It should not be confused with the closely related Kuvi, which the Kandha people also speak. The search found only general reference material and no school materials.
\section{Primary Educational Resources}
None were found in search results. No primers, textbooks or school programmes for Kui turned up.
\section{Linguistic and Reference Materials}
\begin{itemize}
\item Wikipedia, Kui language (India), a general overview of the language: \url{https://www.wikipedia.com/wiki/Khondi}
\item Wikipedia, Khond language, a related overview: \url{https://en.wikipedia.org/wiki/Khond_language}
\item Wikipedia, Kuvi language, for the related but distinct language: \url{https://en.wikipedia.org/wiki/Kuvi_language}
\end{itemize}
No Kui dictionary or grammar was found in the search results.
\section{Religious and Community Resources}
None were found in search results.
\section{Digital and Audio Resources}
None were found in search results.
\section{Curriculum Development Context}
Search results describe Kui as vulnerable in the UNESCO Atlas of the World's Languages in Danger. Documentation appears thin, and no published teaching materials were located. Curriculum work would have to start from fieldwork and from library catalogues of older grammars and word lists, which this search did not verify.
\end{document}